\documentclass[12pt,a4paper]{article}
\usepackage[utf8]{inputenc}
\usepackage{mathptmx}
\usepackage[margin=1in]{geometry}
\usepackage{amsmath,amssymb}
\usepackage{graphicx}
\usepackage{float}
\usepackage{setspace}
\usepackage[round]{natbib}
\usepackage[colorlinks=true,linkcolor=blue,citecolor=blue,urlcolor=blue]{hyperref}
\onehalfspacing

\title{\textbf{The Weak Memory of the Geodynamo: A Random Matrix Theory Verdict on Geomagnetic Reversals}\\[0.4em]
\large Reversal Spacings Lie Between Poisson and GOE, and the Strength of Memory Depends on the Shortest Polarity Intervals}
\author{Ruqing Chen\\ GUT Geoservice Inc., Montreal, Canada\\ \texttt{ruqing@hotmail.com}}
\date{June 2026}

\begin{document}
\maketitle

\begin{abstract}
\noindent A random matrix theory (RMT) program has tested whether single-source, charge-and-release systems exhibit Wigner--Dyson level repulsion, finding repulsion in systems with a depletable reservoir and Poisson statistics where no such memory exists. Here we turn the diagnostic on a long-standing open question in geomagnetism: does the geodynamo have memory? Geomagnetic polarity reversals are an ideal discrete, independently dated event sequence---the events are marine-magnetic-anomaly and radioisotopically dated chron boundaries, not features extracted from a continuous curve, so no extraction artifact applies. The paleomagnetic literature, fitting gamma distributions to chron durations, has reached an ambiguous verdict: the shape parameter $k\approx1.2$--$1.6$ is only slightly above the Poisson value $k=1$, leading some to conclude the reversal process is ``essentially free of long-term memory'' while others infer a weak memory from the systematic variation of reversal frequency. The RMT spacing ratio $\langle r\rangle$ is an independent, memory-sensitive diagnostic that sharpens this verdict. Analyzing the 201 reversals of the well-established C-sequence (0--83 Ma, Cande \& Kent 1995), with local unfolding to remove the long-term reversal-frequency trend, we find $\langle r\rangle = 0.457$ (95\% CI [0.43, 0.51]), which sits $3.0\sigma$ \emph{above} the Poisson value (0.386) and $3.3\sigma$ \emph{below} the GOE value (0.531). The geodynamo therefore has genuine but \emph{weak} memory: reversals are not memoryless, but the repulsion is far weaker than in the crustal charge-and-release systems of this program. The result is robust to unfolding (which lowers CV from 1.49 to 0.78 while leaving $\langle r\rangle$ unchanged) and consistent with the literature gamma index. Crucially, $\langle r\rangle$ rises monotonically toward GOE as the shortest polarity intervals are removed (from 0.457 with all chrons to 0.534 with the shortest 25\% removed), so the inferred memory strength depends directly on the reality of the briefest chrons (cryptochrons/microchrons), a quantity the paleomagnetic community itself regards as uncertain. The geodynamo occupies the weak-memory end of the program's spectrum, consistent with its nature as a turbulent, chaotic magnetohydrodynamic system in which any reservoir memory is strongly diluted.

\vspace{0.8em}
\noindent\textbf{Keywords:} random matrix theory; geomagnetic reversals; geodynamo; polarity timescale; level repulsion; cryptochrons; charge-and-release; memory

\vspace{0.4em}
\noindent\textbf{Code and data:} \url{https://github.com/Ruqing1963/geomagnetic-reversal-rmt}

\vspace{0.2em}
\noindent\textbf{Companion studies:} stratigraphy \url{https://zenodo.org/records/20774581}; seismotectonics \url{https://zenodo.org/records/20768130}; mantle plumes \url{https://zenodo.org/records/20768420}; metallogeny \url{https://zenodo.org/records/20768849}; evolution \url{https://zenodo.org/records/20783763}; hydrogeology \url{https://zenodo.org/records/20780389}; solar flares \url{https://zenodo.org/records/20784967}; orbital architecture \url{https://zenodo.org/records/20785613}; lunar craters \url{https://zenodo.org/records/20787412}
\end{abstract}

\section{Introduction}

A random matrix theory (RMT) program has tested a single hypothesis across the Earth and stellar sciences: that the spacing of events from a single, long-memory charge-and-release source exhibits Wigner--Dyson level repulsion, while memoryless or superposed processes give Poisson statistics \citep{chen2026plume,chen2026metal,chen2026solar,chen2026lunar}. The program has accumulated both positive results (mantle plumes, ore systems, geysers near GOE) and instructive negatives (a single solar active region is Poisson; lunar craters are spatially random or clustered, never repulsive). The diagnostic's value lies in its even-handedness: it returns repulsion only where a depletable reservoir enforces a recharge interval, and randomness where no such memory exists.

Here we apply the diagnostic to a genuinely open question in geomagnetism: \emph{does the geodynamo have memory?} The Earth's magnetic field, generated by convection in the liquid iron outer core, reverses polarity at irregular intervals of $10^5$ to $10^7$ years. A reversal might be viewed as a charge-and-release event: the dynamo loses dipole stability, reverses, and must rebuild a stable configuration, which would impose a minimum recovery time and hence level repulsion in the spacing of reversals. Alternatively, reversals might be triggered randomly by the chaotic turbulence of the core, with no memory between events, giving Poisson statistics. The paleomagnetic community has long debated this. Fitting gamma distributions to polarity-chron durations yields a shape parameter $k\approx1.2$--$1.6$, only slightly above the memoryless Poisson value $k=1$; \citet{lowrie2004} conclude the process is ``essentially free of long-term memory,'' while the pronounced secular variation of reversal frequency (near-zero during the Cretaceous superchron, rising to $\sim$4 Myr$^{-1}$ in the Cenozoic) has been read by others as evidence of long-term core memory. The verdict has remained ambiguous because the gamma index is a weak discriminator near $k=1$.

The RMT spacing ratio $\langle r\rangle$ offers an independent, more memory-sensitive diagnostic. It measures the local rigidity of the sequence directly from consecutive interval ratios, is parameter-free, and distinguishes Poisson (0.386), GOE (0.531), and GUE (0.603) cleanly. We show that it places the geodynamo unambiguously between Poisson and GOE---weak but real memory---and, more usefully, that it pinpoints exactly what the answer depends on: the reality of the shortest polarity intervals.

\section{Data and Methods}

\subsection{The reversal sequence}

We use the C-sequence geomagnetic polarity timescale of \citet{candekent1995} (CK95), the standard, well-calibrated record of reversals from 0 to 83 Ma, anchored by marine magnetic anomalies, radioisotopic ($^{40}$Ar/$^{39}$Ar) ages on selected reversals, and astronomical tuning. This interval comprises 201 reversals (200 polarity chrons) with a mean chron duration of 0.41 Myr. We deliberately restrict the analysis to this single, well-established, stationary-after-detrending interval. We do \emph{not} include the M-sequence (Late Jurassic--Early Cretaceous): the absolute ages of the oldest reversals are less certain, and we verified that a naive, evenly-spaced rendering of M-sequence boundaries produces a spuriously high $\langle r\rangle$, a data artifact we decline to report. We treat the $\sim$38-Myr Cretaceous Normal Superchron (83--121 Ma) as a known, separate anomaly rather than as a single interval in the analysis.

Crucially, reversals are \emph{discrete, independently dated events}, not peaks extracted from a continuous proxy curve. The pipeline therefore carries none of the extraction artifacts that afflict continuous records; the sequence is directly analogous to the discrete, dated event sequences (earthquakes, ore-vein ages, extinctions) analyzed elsewhere in the program.

We emphasize one further point to prevent a cross-reading with the spatial branch of this program. Although polarity reversals are originally read from marine magnetic anomaly stripes on the seafloor (a spatial pattern), the CK95 timescale we use has already converted those stripe widths into event \emph{ages} via independent calibration: the reversal sequence is an explicitly one-dimensional \emph{temporal} point process (event ages in Ma). The temporal Poisson baseline of $0.386$ is therefore mathematically exact here, distinct from the two-dimensional spatial null ($\langle r\rangle\approx0.57$) required for unconstrained surface point fields such as lunar craters \citep{chen2026lunar}. The correct null is set by the dimensionality of the analyzed quantity, and here that quantity is one-dimensional time.

\subsection{Unfolding and diagnostics}

Reversal frequency varies systematically over the C-sequence (slower toward the Cretaceous, faster in the Cenozoic). This large-scale non-stationarity inflates the raw coefficient of variation (CV~$=1.49$) and must be removed before the local spacing statistics are meaningful---exactly the role played by unfolding throughout this program. We apply local unfolding: the cumulative reversal count $N(t)$ is fit with a smooth spline, and spacings are rescaled by the local mean density. After unfolding, CV falls to 0.78, confirming that the frequency trend, not the local rhythm, dominated the raw variance.

For the unfolded spacings we compute the spacing ratio $\langle r\rangle$, the coefficient of variation, the Brody parameter, the gamma shape index $k$ (the literature's diagnostic, for comparison), KS distances to Poisson/GOE/GUE, and 95\% bootstrap confidence intervals. We assess significance against both a Poisson null and a GOE null generated with matched $n$ and span and passed through the identical unfolding.

\section{Results}

\subsection{Weak memory: between Poisson and GOE}

The 201 C-sequence reversals, after local unfolding, give a spacing ratio
\begin{equation}
\langle r\rangle = 0.457 \pm 0.019 \quad (95\%\ \mathrm{CI}\ [0.43,\,0.51]),
\end{equation}
with CV~$=0.78$, Brody $\beta=0.33$, and gamma index $k=0.91$ (Figure~\ref{fig:geodynamo}a, Table~\ref{tab:main}). This value sits well inside the gap between the memoryless and repulsive references. Against a Poisson null (mean $\langle r\rangle=0.389$, generated with matched $n$ and span), the observed value is $3.0\sigma$ \emph{above}: only 0.2\% of Poisson realizations reach 0.457. Against a GOE null (mean 0.531), it is $3.3\sigma$ \emph{below}. Both endpoints are rejected at high confidence. The geodynamo is therefore neither memoryless nor strongly repulsive: it has \emph{weak memory}.

This is a sharper statement than the gamma-fitting literature could make. Our gamma index ($k=0.91$ on all chrons, rising to 1.14 when the shortest chrons are removed) reproduces the published range $k\approx1.2$--$1.6$ and confirms the method, but the gamma fit near $k=1$ is a weak discriminator. The spacing ratio resolves the same data into a clear position: $3\sigma$ above Poisson, $3\sigma$ below GOE.

\begin{table}[H]
\centering
\caption{RMT diagnostics for the C-sequence geomagnetic reversal record (0--83 Ma, 201 reversals). The geodynamo sits between Poisson and GOE; the inferred memory strengthens as the shortest polarity intervals are excluded.}
\label{tab:main}
\small
\begin{tabular}{lccccl}
\toprule
Analysis & $n$ & $\langle r\rangle$ & CV & gamma $k$ & position \\
\midrule
Raw (no unfolding) & 201 & 0.456 & 1.49 & 0.91 & Poisson--GOE \\
\textbf{Unfolded (frequency trend removed)} & \textbf{201} & \textbf{0.457} & \textbf{0.78} & \textbf{0.91} & \textbf{weak memory} \\
Drop shortest 10\% chrons & 181 & 0.487 & 0.73 & 1.14 & toward GOE \\
Drop shortest 20\% chrons & 161 & 0.530 & 0.67 & --- & at GOE \\
\midrule
Poisson null (matched) & 201 & 0.389 & --- & 1.0 & reference \\
GOE null (matched) & 201 & 0.531 & --- & --- & reference \\
\bottomrule
\end{tabular}
\end{table}

\subsection{The cryptochron dependence}

The most informative result is how $\langle r\rangle$ responds to the shortest polarity intervals (Figure~\ref{fig:geodynamo}b). Removing the shortest chrons raises $\langle r\rangle$ monotonically: 0.457 (all chrons) $\to$ 0.487 (shortest 10\% removed) $\to$ 0.530 (shortest 20\%) $\to$ 0.534 (shortest 25\%). The very brief polarity intervals---cryptochrons and microchrons, some lasting only a few thousand years---are precisely what pull the statistic toward Poisson. If these brief intervals are all real reversals, the geodynamo has only weak memory ($\langle r\rangle=0.457$); if a substantial fraction are recording artifacts or non-global excursions rather than true global reversals, the underlying reversal process is more rigid, approaching GOE.

This is not a statistical nuisance but a direct connection to an active controversy in paleomagnetism. The reality, globality, and even the definition of the briefest events is genuinely contested: \citet{laj2007} distinguish ``excursions'' (local departures that do not reach fully reversed directions) from ``microchrons'' (brief but full polarity intervals), and well-studied cases such as the Mono Lake--Laschamp feature(s) remain debated as to whether they represent one event or two, and whether they are globally recorded. The RMT diagnostic translates this qualitative debate into a quantitative lever: the memory strength of the geodynamo, as measured by $\langle r\rangle$, is a direct function of how many of the shortest intervals are accepted as true reversals. A future, cryptochron-cleaned polarity timescale would move $\langle r\rangle$ to a definite value along the 0.46--0.53 range we map here.

\begin{figure}[H]
\centering
\includegraphics[width=\textwidth]{figs/fig1_geodynamo.pdf}
\caption{(a) Unfolded reversal-spacing distribution for the C-sequence (0--83 Ma, 201 reversals). The data (blue) lie between Poisson (dotted) and GOE (dashed), at $\langle r\rangle=0.457$: $3.0\sigma$ above Poisson, $3.3\sigma$ below GOE---weak but real memory. (b) The cryptochron dependence: $\langle r\rangle$ rises monotonically toward GOE as the shortest polarity intervals are removed, so the inferred memory strength depends on the reality of the briefest chrons.}
\label{fig:geodynamo}
\end{figure}

\section{Discussion}

\subsection{The geodynamo on the memory spectrum}

The program has, across its systems, traced a spectrum of memory strength. At the strong-memory (GOE) end are the crustal charge-and-release systems: an ore-forming hydrothermal cell that cannot re-precipitate in an already-sealed vein, a volcanic or geyser reservoir that must refill, a mantle plume conduit. At the memoryless (Poisson) end are systems with no single depletable reservoir: a single solar active region near its depletion limit, or memoryless impact bombardment. The geodynamo, at $\langle r\rangle=0.457$, falls firmly in between, near the weak-memory end.

This placement is physically coherent. The outer core is a vigorously turbulent, chaotic magnetohydrodynamic system. Simple dynamo models---even two-disk systems---produce chaotic polarity switching, and the real core is far more complex. Whatever ``reservoir'' a reversal depletes (dipole stability, a particular flow configuration), the recovery is governed by strongly chaotic dynamics that dilute any clean recharge interval. One would therefore expect the geodynamo to retain only weak memory: enough that a reversal is not instantly followed by another (hence $\langle r\rangle$ above Poisson), but far from the rigid, near-deterministic recharge of a sealed crustal reservoir (hence far below GOE). The RMT result quantifies this intuition: core memory is real but heavily diluted by turbulence.

The charge-and-release cycle can be made concrete in magnetohydrodynamic (MHD) terms. The ``release'' is the collapse of the axial dipole into a transient multipolar state during a reversal; the ``recharge'' is the regeneration of a stable axial dipole, organized by convection and the Coriolis force in the rotating core. The reason the memory is weak rather than strong lies in the very high magnetic Reynolds number of the outer core, where advection dominates magnetic diffusion. Turbulent eddies continuously generate localized reverse-flux patches at the core--mantle boundary, and such a patch can grow and trigger a premature, sometimes aborted reversal before the global dipole has fully re-stabilized. These premature events are exactly the brief intervals---cryptochrons and microchrons---that our analysis shows pull $\langle r\rangle$ toward Poisson. The cryptochron dependence of Section~3.2 is therefore not merely a data uncertainty but the statistical fingerprint of this physical dilution: turbulent reverse-flux patches inject short, quasi-random intervals into what would otherwise be a more rigid recharge rhythm, lowering the Wigner--Dyson memory toward the memoryless limit.

\subsection{Why the spacing ratio outperforms the gamma fit}

The paleomagnetic literature's gamma index has hovered ambiguously near $k=1$ for decades, supporting both ``no memory'' and ``some memory'' readings. The spacing ratio is a better instrument here for two reasons. First, it is local: it compares each interval only to its immediate neighbors, so it is insensitive to the slow frequency drift that a global distribution fit must contend with (we remove that drift explicitly by unfolding). Second, it is calibrated against the same universal references (Poisson/GOE/GUE) used throughout physics, giving an absolute scale of memory rather than a single shape parameter whose interpretation near $k=1$ is contested. The convergence of our gamma index with the literature confirms we analyze the same data; the spacing ratio simply extracts a clearer verdict from it.

\subsection{A falsifiable handle on cryptochrons}

The cryptochron dependence turns an abstract debate into a measurement. Because $\langle r\rangle$ moves predictably as brief intervals are included or excluded, the diagnostic provides a quantitative target: as the global reality of the shortest events is settled by future high-resolution magnetostratigraphy, the geodynamo's true memory strength is pinned to a definite point on the 0.46--0.53 interval. If most brief events prove to be genuine global reversals, the dynamo's memory is weak ($\sim$0.46); if many are local excursions or artifacts, the underlying process is moderately repulsive ($\sim$0.53). This is a concrete, testable prediction that links a statistical diagnostic to an observational program.

\subsection{Limitations}

\emph{Single interval.} We restrict to the CK95 C-sequence (0--83 Ma) to ensure a well-calibrated, stationary-after-detrending record; we do not extend to the M-sequence, whose less certain ages we found could produce artifacts. \emph{Superchron excluded.} The Cretaceous Normal Superchron is treated as a separate known anomaly, not folded into the spacing statistics---not as an extreme long-waiting-time draw from the standard reversal distribution, but because the leading hypothesis holds that superchrons represent a fundamentally distinct geodynamo regime driven by altered boundary conditions (notably reduced core--mantle boundary heat flux) rather than a stochastic tail of the normal reversal process. Its origin remains debated, but on either reading it should not be pooled with the standard-regime spacings analyzed here. \emph{Cryptochron census.} Our central value (0.457) uses all chrons in CK95; the true memory strength is bracketed, not pinned, until the global reality of the briefest intervals is resolved---which is precisely the dependence we quantify. \emph{Reversal vs.\ chron.} We analyze chron-boundary spacings; alternative event definitions (e.g., midpoints of transitional intervals) would shift absolute values slightly but not the qualitative weak-memory placement.

\section{Conclusions}

Using the 201 reversals of the CK95 C-sequence polarity timescale (0--83 Ma), with local unfolding to remove the reversal-frequency trend, we find that geomagnetic reversals have a spacing ratio $\langle r\rangle=0.457$, lying $3.0\sigma$ above the Poisson value and $3.3\sigma$ below GOE. The geodynamo has genuine but weak memory: reversals are not memoryless, but the repulsion is far weaker than in the crustal charge-and-release systems of this program. The inferred memory strength depends directly on the reality of the shortest polarity intervals, rising toward GOE as cryptochrons are excluded---turning a long-standing qualitative debate into a quantitative, falsifiable measurement. The geodynamo occupies the weak-memory end of the program's spectrum, exactly as expected for a turbulent, chaotic magnetohydrodynamic system in which reservoir memory is strongly diluted. This is the tenth system in the program and the first to place a deep-Earth, core-generated process on the memory spectrum, between the strong repulsion of crustal reservoirs and the randomness of memoryless bombardment.

\section*{Code and Data Availability}
All code, the reversal-age table, and analysis outputs are at \url{https://github.com/Ruqing1963/geomagnetic-reversal-rmt}. Reversal ages are from the Cande \& Kent (1995) C-sequence GPTS. The pipeline reproduces every number and figure.

\begin{thebibliography}{99}
\small
\bibitem[Cande \& Kent(1995)]{candekent1995} Cande, S. C., \& Kent, D. V. (1995). Revised calibration of the geomagnetic polarity timescale for the Late Cretaceous and Cenozoic. \emph{Journal of Geophysical Research}, 100, 6093--6095.
\bibitem[Lowrie \& Kent(2004)]{lowrie2004} Lowrie, W., \& Kent, D. V. (2004). Geomagnetic polarity timescales and reversal frequency regimes. In \emph{Timescales of the Paleomagnetic Field}, AGU Geophysical Monograph 145, 117--129.
\bibitem[Laj \& Channell(2007)]{laj2007} Laj, C., \& Channell, J. E. T. (2007). Geomagnetic excursions. In \emph{Treatise on Geophysics}, vol. 5, 373--416.
\bibitem[Tauxe(2010)]{tauxe2010} Tauxe, L. (2010). \emph{Essentials of Paleomagnetism}. University of California Press.
\bibitem[Chen(2026b)]{chen2026plume} Chen, R. (2026b). Level repulsion in the deep Earth: mantle-plume eruption rhythms. Zenodo, \url{https://zenodo.org/records/20768420}.
\bibitem[Chen(2026c)]{chen2026metal} Chen, R. (2026c). Charge-and-release rhythms in ore-forming systems. Zenodo, \url{https://zenodo.org/records/20768849}.
\bibitem[Chen(2026g)]{chen2026solar} Chen, R. (2026g). The depletion limit of level repulsion: solar flares from a single active region. Zenodo, \url{https://zenodo.org/records/20784967}.
\bibitem[Chen(2026i)]{chen2026lunar} Chen, R. (2026i). The passive baseline: why lunar impact cratering is not charge-and-release. Zenodo, \url{https://zenodo.org/records/20787412}.
\bibitem[Bohigas et al.(1984)]{bohigas1984} Bohigas, O., Giannoni, M.-J., \& Schmit, C. (1984). Characterization of chaotic quantum spectra and universality of level fluctuation laws. \emph{Physical Review Letters}, 52, 1.
\end{thebibliography}

\end{document}
